% ============================================================
\section{Cálculos de diseño complementarios}\label{anx:calculos}
% ============================================================

\subsection{Verificación térmica del regulador KIA7805}
La corriente del riel de \SI{5}{\volt} en el peor caso (pinza haciendo fuerza) se estima con
el incremento de consumo medido y valores típicos para el resto de las cargas:
\begin{equation}
  I_{5V} \approx \underbrace{0{,}54}_{\text{servo}} + \underbrace{0{,}03}_{\text{LCD}} + \underbrace{0{,}05}_{\text{riel \SI{3,3}{\volt}}} \approx \SI{0,62}{\ampere}.
\end{equation}
La potencia disipada por el regulador lineal es
\begin{equation}
  P_D = (V_I - V_O)\,I_O = (11{,}2 - 5)\cdot 0{,}62 \approx \SI{3,8}{\watt}.
\end{equation}
Con resistencias térmicas típicas del encapsulado TO-220 ($R_{\theta JC}\approx\SI{5}{\celsius\per\watt}$,
$R_{\theta JA}\approx\SI{65}{\celsius\per\watt}$ sin disipador), el aumento de temperatura sin disipador
sería de $3{,}8\cdot 65\approx\SI{250}{\celsius}$, que activaría la protección térmica
interna. Para $T_J\le\SI{125}{\celsius}$ con $T_A=\SI{40}{\celsius}$ se necesita
\begin{equation}
  R_{\theta SA} \le \frac{T_J - T_A}{P_D} - R_{\theta JC} - R_{\theta CS}
  = \frac{125-40}{3{,}8} - 5 - 0{,}5 \approx \SI{17}{\celsius\per\watt}.
\end{equation}
En reposo ($I_{5V}\approx\SI{0,1}{\ampere}$, $P_D\approx\SI{0,6}{\watt}$) el regulador funciona
sin problemas aun sin disipador. El caso crítico es mantener la pinza cerrada con fuerza
durante mucho tiempo. En el prototipo el KIA7805 \textbf{lleva un disipador} de aletas para
encapsulado TO-220; de acuerdo con el cálculo anterior basta con que su resistencia térmica sea
$R_{\theta SA}\lesssim\SI{17}{\celsius\per\watt}$ —un valor holgado para un disipador pequeño con
clip—, de modo que la unión se mantiene por debajo de \SI{125}{\celsius} aun en el peor caso.
De todos modos conviene verificarlo midiendo la temperatura de la lengüeta con la pinza cerrada
en régimen; a futuro, reemplazar el regulador lineal por uno conmutado (Sección~\ref{sec:conclusiones})
elimina el problema de raíz.

\subsection{Regulador LM317 de \SI{3,3}{\volt}}
\begin{itemize}
  \item \textbf{Ajuste:} $V_O = 1{,}25\,(1+R_2/R_1) + I_{\mathrm{ADJ}}R_2$. Con $R_1=\SI{330}{\ohm}$,
        $R_2\approx\SI{541}{\ohm}$ en el preset de \SI{1}{\kilo\ohm}. El término
        $I_{\mathrm{ADJ}}R_2\approx \SI{50}{\micro\ampere}\cdot\SI{541}{\ohm}\approx\SI{27}{\milli\volt}$
        es despreciable.
  \item \textbf{Carga mínima:} la corriente por $R_1$ es $1{,}25/330\approx\SI{3,8}{\milli\ampere}$,
        dentro de la carga mínima que requiere el LM317 para regular.
  \item \textbf{Tensión diferencial:} $V_I - V_O = 5 - 3{,}3 = \SI{1,7}{\volt}$, cercana a la
        tensión de \emph{dropout} del LM317 (aprox.\ \SIrange{1,5}{2}{\volt} según la corriente).
        Con la tolerancia del 7805 (\SI{4,8}{\volt}) el margen es mínimo. Funciona por la baja
        corriente del riel, pero para una próxima revisión conviene un LDO de \SI{3,3}{\volt}.
\end{itemize}

\subsection{Indicador de final de carrera (BC558)}
Si se tiene en cuenta la corriente de base, la corriente de emisor en conducción es
\begin{equation}
  I_E = \frac{V_{CC} - V_F - V_{EB} - V_{B}}{R_E},
  \qquad V_B \approx V_{OL,\mathrm{sensor}} \approx 0
  \;\Rightarrow\; I_E \approx \frac{3{,}3-2{,}0-0{,}7}{220} \approx \SI{2,7}{\milli\ampere}.
\end{equation}
La corriente que absorbe la salida del sensor es la de base
($I_E/(\beta+1)\approx\SI{15}{\micro\ampere}$) más la del \emph{pull-up} de \SI{47}{\kilo\ohm}
($3{,}3/47\,\mathrm{k}\approx\SI{70}{\micro\ampere}$) y la del \emph{pull-up} interno del STM32
(aprox.\ \SI{80}{\micro\ampere}). En total son menos de \SI{0,2}{\milli\ampere}, compatibles con
cualquier salida de colector abierto o con un contacto mecánico.

\subsection{Presupuesto de tiempo de la ISR de TIM2}
Una \emph{rutina de servicio de interrupción} (ISR, \emph{Interrupt Service Routine}) es la
función que el procesador ejecuta de inmediato cuando se dispara un evento de hardware —aquí, el
desborde periódico de TIM2—, suspendiendo el lazo principal hasta que termina. Por eso su
duración debe ser mucho menor que el intervalo entre disparos: si la ISR tardara más que los
\SI{50}{\micro\second} que separan dos desbordes, no quedaría tiempo para atender el USB ni el
resto de las tareas.

La ISR de TIM2 se ejecuta cada $T=\SI{50}{\micro\second}$, es decir, cada 3600 ciclos de CPU a
\SI{72}{\mega\hertz}. En cada llamada recorre los tres motores con lecturas de GPIO (sensores),
comparaciones y, cuando corresponde, escrituras de STEP/DIR y el recálculo de la rampa con una
división entera.

\paragraph{Método de medición propuesto.} La forma práctica de medir su duración es
\emph{instrumentar} el código: poner en alto un pin de prueba (un GPIO libre) en la primera línea
de la ISR y ponerlo en bajo en la última, y observar ese pin con un osciloscopio. El ancho del
pulso es el tiempo de ejecución y el período entre pulsos confirma los \SI{50}{\micro\second}; la
carga de CPU es (ancho del pulso)\,/\,\SI{50}{\micro\second}. Debe medirse el \textbf{peor caso},
que ocurre cuando los tres ejes generan un flanco de paso en el mismo \emph{tick} (máximo de
escrituras de GPIO y tres llamadas a \texttt{CalculateSpeed}, cada una con su división entera).
La medición se hace con los tres ejes moviéndose a la vez —no eje por eje—, porque la ISR siempre
recorre los tres motores en cada llamada. Se estima una carga menor al 15~\%, lo que dejaría
amplio margen para el USB y el lazo principal; la medición con osciloscopio queda pendiente de
realizar.
